\documentclass[10pt,a4paper]{amsart}
\usepackage[margin=19mm]{geometry}
\usepackage{amsmath,amssymb,mathtools}
\usepackage[scr=rsfs,cal=euler]{mathalfa}
\usepackage{xfrac}
\usepackage[unicode,hidelinks,bookmarks=false]{hyperref}
\newcommand{\ie}{i.e.\ }
\newcommand{\cf}{cf.\ }
\newcommand{\resp}{respectively\ }
\newcommand{\Subsection}{\subsection}
\providecommand{\nobreakdash}{\nobreak\hbox{-}}
\setlength{\emergencystretch}{2em}
\multlinegap0pt
\usepackage{amssymb}
\usepackage{xcolor}
\newtheorem{lemma}{Lemma}
\title{Kronecker coefficients via the Giambelli identity for Schur functions}
\author{John M. Campbell}
\date{Source: arXiv:2604.23286v1 (CC BY 4.0)}
\begin{document}
\maketitle

\noindent\emph{Source and licence.} Kronecker coefficients via the Giambelli identity for Schur functions. Version arXiv:2604.23286v1 (CC BY 4.0), distributed by arXiv under the Creative Commons Attribution 4.0 International licence, http://creativecommons.org/licenses/by/4.0/, as recorded in the arXiv licence metadata for this exact version and shown on the abstract page https://arxiv.org/abs/2604.23286v1. Full licence notice: \texttt{licenses/arxiv-2604.23286-CC-BY-4.0.txt}.\par\smallskip

\noindent\textit{Editorial scope.} Counted unit: the article's lemma Lemsize1 (source0.tex lines 1282--1296). The article's complete original proof of it is delivered with it (source0.tex lines 1298--1346, one \texttt{proof} environment of 49 lines). No mathematics is written, repaired or re-derived: the statement and the proof are the article's own lines, in the article's own notation, and no retained line is substituted or re-flowed.  Retained uncounted as the source-local context the proof needs: lines 502--504.  Nothing was omitted from any retained span, so the package carries no omission record.  No retained line contains a citation, so the wrapper carries no bibliography and no bibliography entry.  No retained line contains a figure environment, an image-inclusion command or drawing code, so no image is delivered and the package declares no asset dependency.  Editorial formatting only, in addition: an AMS \texttt{amsart} preamble that restates the article's own declarations and macros used by the retained lines, namely the environments \texttt{lemma} on separate counters, together with the standard packages those lines need.  The retained environments print in this document as Lemma 1 in order of appearance, whereas the article prints the same items with the numbering of its own full document.  The complete native source of the article as distributed in the arXiv e-print for this exact version is bundled unchanged as \texttt{sources/199150/source0.tex}.
\begin{equation}\label{gsignrep}
 g_{(1^{k}) \, \lambda \, \mu} = \text{{\bf 1}}_{\lambda = \mu^{\operatorname{t}}}. 
\end{equation} 

\begin{lemma}\label{Lemsize1}
 Let $a \geq 2$, $b = 2$, $c \geq 1$, $d$, $e$ and $\mathcal{S}$ be integers and let $\nu$ be a partition of $n$, with $a+b+c = d + e = 
 n$ and $d \geq e$ (and we again let $N = n - b + 1$ and $M = n - a$ be as above). Moreover, let $1 \leq \mathcal{S} \leq \lfloor 
 \frac{c+2}{2} \rfloor$, and suppose that 
\begin{enumerate}

\item $\nu = (a+2, 2^{\mathcal{S}-1}, 1^{c+2-2\mathcal{S}})$; and 

\item $\max(a + \mathcal{S}, c + 2 - \mathcal{S}) \leq d \leq a + c + 2 -\mathcal{S}$.

\end{enumerate}

\noindent Then $\left| \mathcal{J}_{d}^{-}(\nu) \right| = 1$, and, moreover, the unique term (in this case) within
 $ \operatorname{triple}_{(d, e) \, (a, b, 1^c) \, \nu}^{(4)}$ is equal to $1$. 
\end{lemma}

\begin{proof}
 We now isolate a family for which the negative contribution 
\(\operatorname{triple}^{(4)}\) collapses to a single term of value \(1\). Let 
 $$ \delta^{\ast} := \big( 2^{\mathcal{S}}, 1^{c+2-2\mathcal{S}} \big) \ \ \ \text{and} 
 \ \ \ M = n - a = c + 2. $$
 We claim that 
\begin{equation}\label{Jdminussingle}
 \mathcal{J}_{d}^{-}(\nu) = \big\{ (\delta^{\ast}, 0, \mathcal{S} ) \big\} 
\end{equation}
 and that the unique term in $\operatorname{triple}^{(4)}$ corresponding to the singleton index set in \eqref{Jdminussingle} is 
 equal to $1$. 

 Let $(\delta, i, r) \in \mathcal{J}_{d}^{-}(\nu)$. From the containment $\mathcal{J}_{d}^{-}(\nu) \subseteq \mathcal{I}^{-}(\nu)$, we obtain 
 the inequality $$ \mathcal{N}_{\nu \, \delta \, i \, a - i} > \mathcal{N}_{\nu \, \delta \, i - 1 \, a - i + 1}, $$ so that $$ \mathcal{N}_{\nu 
 \, \delta \, i \, a - i} > 0. $$ From the definition of our $\mathcal{N}$-sets, there exists a partition $\kappa$ such that $$ \delta \subseteq 
 \kappa \subseteq \nu, $$ and hence 
\begin{equation}\label{deltacontained}
 \delta \subseteq (a+2, 2^{\mathcal{S}-1}, 1^{c+2-2\mathcal{S}}).
\end{equation}
 Since $|\delta| = M = c + 2$, we can conclude, from \eqref{deltacontained}, that the number of boxes of $\nu$ below the first row is
 $2 (\mathcal{S}-1) + (c + 2 - 2 \mathcal{S}) =c$. So, a given subpartition of $\nu$ of size $c+2$ necessarily consists of $c$ boxes
 under the initial row and $2$ boxes in the initial row. The unique such partition is $$ \delta = (2^\mathcal{S}, 1^{c + 2 - 
 2\mathcal{S}}) = \delta^{\ast}, $$ i.e., so that every element in $\mathcal{J}_{d}^{-}(\nu)$ is of the form $(\delta^{\ast}, i, r)$.
 
 Writing $\delta = \delta^{\ast}$, we see that $ \nu / \delta^{\ast} $ consists of $a$ boxes, and all of these boxes are in the initial row. 
 So, for $i$ such that $0 \leq i \leq a$, there exists a unique intermediate partition $$ \kappa^{(i)} = \big( i + 2, 2^{\mathcal{S}-1}, 
 1^{c + 2 - 2 \mathcal{S}} \big) $$ such that $\kappa^{(i)} / \delta^{\ast}$ is a horizontal strip of size $i$ and $\nu / \kappa^{(i)}$ is a 
 horizontal strip of size $a - i$. We thus obtain that $$ \mathcal{N}_{\nu \, \delta^{\ast} \, i \, a - i} = 1 $$ for $i$ such that $0 \leq i \leq 
 a$. Also, if $i \geq 1$, then $$ \mathcal{N}_{\nu \, \delta^{\ast} \, i - 1 \, a - i + 1 } = 1.$$ From the requirement that $$ \mathcal{N}_{\nu 
 \, \delta^{\ast} \, i \, a - i } > \mathcal{N}_{\nu \, \delta^{\ast} \, i - 1 \, a - i +1}, $$ we find that $ i = 0$. 

 Since $b =2$, we find that $$ g_{(M-r,r) \, (b-1 \, 1^{c+1}) \, \delta^{\ast} } = g_{(M-r, r) \, (1^{M}) \, \delta^{\ast}}. $$ The Kronecker 
 coefficient identity in \eqref{gsignrep} associated with sign representations allows us to conclude that $$ g_{(M-r,r) \, (1^M) \, 
 \delta^{\ast} } = 1 \Longleftrightarrow \delta^{\ast} = (M-r,r)^{\operatorname{t}}. $$ Since $$ (M-\mathcal{S}, 
 \mathcal{S})^t=(2^{\mathcal{S}},1^{M-2\mathcal{S}})=\delta^\ast, $$ we can conclude that the biconditional equivalence 
 $$ \delta^\ast=(M-r,r)^t \Longleftrightarrow r=\mathcal S $$ holds. Consequently, the unique element in 
 the index set $\mathcal{J}_{d}^{-}(\nu)$ is $(\delta^{\ast}, 0, \mathcal{S})$. 

 So, it remains to evaluate the unique term in $\operatorname{triple}^{(4)}_{(d,e)\,(a,b,1^c)\,\nu}$ corresponding to $(\delta^\ast, 0, 
 \mathcal S)$. Since $i=0$, this term is $$ c_{(a,0) \, \delta^\ast}^{\nu} \, g_{(M-\mathcal S,\mathcal S)\,(1^M)\,\delta^\ast}, $$ with 
 the reduction $(a,0)=(a)$. Since $\nu/\delta^\ast$ consists of $a$ boxes, with each such box being in the initial row, the Pieri rule gives 
 us that $ c_{(a)\,\delta^\ast}^{\nu}=1$. Moreover, from the sign representation identity in \eqref{gsignrep}, the equality 
 $$ g_{ (M - \mathcal{S}, \mathcal{S}) \, (1^{M}) \, \delta^{\ast} } = \text{{\bf 1}}_{ (M - \mathcal{S}, \mathcal{S} )^{\operatorname{t}} 
 = \delta^{\ast}} $$ holds true. Since $$ (M - \mathcal{S}, \mathcal{S})^{\operatorname{t}} = (2^{\mathcal{S}}, 1^{M - 2 \mathcal{S}}) = 
 (2^{\mathcal{S}}, 1^{c + 2 - 2\mathcal{S}}) = \delta^\ast, $$ we may obtain that $ g_{(M - \mathcal{S}, \mathcal{S}) \, (1^M) \, 
 \delta^\ast}=1$. As a consequence, the unique term is $1\cdot 1=1$, and hence the equality 
 $ |\mathcal J_d^{-}(\nu)|=1$. So, the unique term in $\operatorname{triple}^{(4)}_{(d,e)\,(a,b,1^c)\,\nu}$
 is equal to $1$, as required. 
\end{proof}

\end{document}
